\begin{theorem}[Elementary-graph criterion]
Let $G=(V,E)$ satisfy the paper's criticality convention, here expressed by
the conjunction that deleting every edge lowers $\chi'(G)$ and that every
vertex is incident with an edge.  Suppose that
$\chi'(G)=k+1$ and $k\geq\Delta(G)+1$.  Then the following four conditions
are equivalent:
\begin{enumerate}
\item[(i)] $G$ is elementary.
\item[(ii)] For every edge identity $e\in E$ and every proper partial
$k$-edge-coloring $\varphi$ which leaves exactly $e$ uncolored, the whole
vertex set $V$ is elementary with respect to $\varphi$.
\item[(iii)] There are an edge identity $e\in E$ and such a coloring
$\varphi$ for which $V$ is elementary with respect to $\varphi$.
\item[(iv)] There are an edge identity $e\in E$, such a coloring $\varphi$,
and a set $X\subseteq V$ containing both ends of $e$, such that $X$ is both
elementary and strongly closed with respect to $\varphi$.
\end{enumerate}
Equivalence means the conjunction
$(\mathrm{i}\leftrightarrow\mathrm{ii})\land
(\mathrm{ii}\leftrightarrow\mathrm{iii})\land
(\mathrm{iii}\leftrightarrow\mathrm{iv})$.
\end{theorem}
\begin{proof}
We use the definitions of criticality, chromatic index, density, and
elementariness from \noderef{Critical}, \noderef{ChromaticIndex},
\noderef{Density}, \noderef{DensityCeil}, and \noderef{Elementary}.  Missing
colors and the three colored-set properties have the meanings fixed in
\noderef{MissingColors}, \noderef{ColoredVertexSetElementary},
\noderef{ColoredVertexSetClosed}, and
\noderef{ColoredVertexSetStronglyClosed}.  Properness and the assertion that
exactly $e$ is uncolored are those of \noderef{IsProperEdgeColoring} and
\noderef{IsOneEdgeUncolored}.  We prove the implications cyclically.

Assume (i).  Since $G$ is elementary and $\chi'(G)=k+1$, we have
$\lceil\Gamma(G)\rceil=k+1$, and hence $\Gamma(G)>k$.  By the finite maximum
in \noderef{Density}, there is an odd set $U$ of cardinality at least three
with
\[
  2|E_G(U)|>k(|U|-1).                                      \tag{1}
\]
Every edge of $G$ has both ends in $U$.  Indeed, if an edge identity $f$
were not internal to $U$, criticality would give
$\chi'(G-f)<k+1$, hence a proper coloring of $G-f$ with at most $k$ colors.
All edges of $E_G(U)$ would remain present.  Each color class is a matching
on the odd set $U$, so it contains at most $(|U|-1)/2$ internal edges.
Summing over at most $k$ colors would contradict (1).  Because every vertex
is incident with an edge, the fact that all edges are internal to $U$ now
implies $U=V$.

Fix an edge $e$ and a proper partial $k$-edge-coloring $\varphi$ leaving
exactly $e$ uncolored.  Put $n=|V|$ and $m=|E|$.  The preceding paragraph
shows that $n$ is odd and (1) says $m>k(n-1)/2$.  Since $m$ is integral,
$m-1\geq k(n-1)/2$.  On the other hand, the $m-1$ colored edges are
partitioned into $k$ color classes, each a matching on the odd set $V$ and
therefore of size at most $(n-1)/2$.  Thus
\[
  m-1=k(n-1)/2,
\]
and equality must hold in the bound for every color class.  Every color
class consequently saturates exactly $n-1$ vertices and misses exactly one
vertex.  No color can therefore be missing at two distinct vertices.  By
\noderef{ColoredVertexSetElementary}, $V$ is elementary with respect to
$\varphi$.  Since $e$ and $\varphi$ were arbitrary, (ii) follows.

Assume (ii).  The equality $\chi'(G)=k+1$ implies that $E$ is nonempty:
an edgeless graph has chromatic index zero.  Choose $e\in E$.  Criticality
gives $\chi'(G-e)<k+1$, and hence $\chi'(G-e)\leq k$.  The existence part of
\noderef{ChromaticIndexSpec}, applied to $G-e$, supplies a proper coloring
with $\chi'(G-e)$ colors; regarding those colors as colors in $[k]$ and
assigning the value ``uncolored'' only to $e$ gives a proper partial
$k$-edge-coloring of $G$ leaving exactly $e$ uncolored.  Condition (ii)
makes $V$ elementary for this choice.  This proves (iii).

Assume (iii), witnessed by $e$ and $\varphi$, and take $X=V$.  Both ends of
$e$ belong to $X$.  The boundary of the whole vertex set is empty, so the
closure condition and the distinct-boundary-color condition are both
vacuous.  Hence $X$ is strongly closed, while its elementariness is the
hypothesis in (iii).  This proves (iv).

Finally assume (iv), witnessed by $e$, $\varphi$, and $X$.  Let $x$ be one
end of $e$.  The two ends are distinct by the loopless law in
\noderef{LooplessMultigraph}.  Since $e$ is the only uncolored edge and is
incident with $x$, at most $d_G(x)-1$ colors are present at $x$.  The bound
$k\geq\Delta(G)+1$ therefore gives at least two colors missing at $x$.
Choose one and call it $\alpha$.

For a color $c$, let $m_c$ be the number of vertices of $X$ missing $c$, let
$b_c$ be the number of boundary-edge identities of $X$ colored $c$, and let
$a_c$ be the number of internal edges of $X$ colored $c$.  Counting
incidences between vertices of $X$ and $c$-colored edges gives
\[
  2a_c+b_c=|X|-m_c.                                       \tag{2}
\]
Indeed, properness says that each vertex not missing $c$ has exactly one
incident $c$-edge; an internal edge contributes two such incidences and a
boundary edge one.  Because $X$ is elementary and $\alpha$ is missing at
$x$, we have $m_\alpha=1$.  Because $X$ is closed, no boundary edge has
color $\alpha$, so $b_\alpha=0$.  Equation (2) says
$2a_\alpha=|X|-1$, and therefore $|X|$ is odd.  The set $X$ contains the
two distinct ends of $e$, so its odd cardinality is at least three.

Let $q$ be the number of colors missing somewhere on $X$.  Elementariness
makes the missing-color sets at distinct vertices disjoint, so $q$ is the
sum of their cardinalities.  Summing, over $X$, the identity ``number of
missing colors equals $k$ minus colored degree'' and again counting internal
incidences twice and boundary incidences once yields
\[
 q=k|X|-2|E_{G-e}(X)|-|\partial_G(X)|.                    \tag{3}
\]
The edge $e$ is internal to $X$, so it does not occur in the boundary and
every boundary edge is colored.  Strong closure says that their colors are
pairwise distinct, while closure says that none of those colors is among
the $q$ colors missing on $X$.  Thus
$|\partial_G(X)|\leq k-q$.  Substitution in (3) gives
\[
  2|E_{G-e}(X)|\geq k(|X|-1).
\]
Adding the internal uncolored edge $e$ gives
\[
  2|E_G(X)|>k(|X|-1).
\]
Since $X$ is an eligible odd set, \noderef{DensityValue} and
\noderef{Density} imply $\Gamma(G)>k$, and hence
$\lceil\Gamma(G)\rceil\geq k+1=\chi'(G)$.  The reverse inequality is
\noderef{DensityLowerBound}.  Equality follows, so $G$ is elementary.
This proves (iv)$\Rightarrow$(i) and completes all three asserted
equivalences.
\end{proof}
